\documentclass[11pt,letterpaper]{article}
\usepackage[T1]{fontenc}
\usepackage{lmodern,amsmath,amssymb,amsthm,microtype}
\usepackage[margin=1in]{geometry}
\usepackage[colorlinks=true,linkcolor=blue,citecolor=blue,urlcolor=blue]{hyperref}
\hypersetup{pdftitle={Counterexamples to Satoh's SL(2) representation-filtration conjecture},
pdfauthor={Alec Kriebel},
pdfsubject={Explicit counterexamples to a filtration equality formulated in 2018},
pdfkeywords={free group automorphisms, Andreadakis filtration, representation algebra, polynomial identity}}
\newtheorem{theorem}{Theorem}
\newtheorem{lemma}[theorem]{Lemma}
\newcommand{\Aut}{\operatorname{Aut}}
\newcommand{\Mat}{\operatorname{Mat}}
\newcommand{\SL}{\operatorname{SL}}
\newcommand{\Q}{\mathbb{Q}}
\newcommand{\Z}{\mathbb{Z}}
\newcommand{\C}{\mathbb{C}}
\title{Counterexamples to Satoh's \texorpdfstring{$\SL(2)$}{SL(2)}\protect\\ representation-filtration conjecture}
\author{Alec Kriebel\\\small Independent researcher\\
\small\href{https://orcid.org/0009-0001-9320-500X}{ORCID: 0009-0001-9320-500X}}
\date{15 October 2026}
\begin{document}
\maketitle
\begin{abstract}
Let $A_n(k)$ be the integral Andreadakis filtration of the automorphism group
of a free group, and let $D_n(k)$ be the kernel of its action on
$J/J^{k+1}$ in the rational $\SL(2,\C)$ matrix-entry algebra at the trivial
representation. We construct elements of $D_n(5)\setminus A_n(5)$ for every
$n\geq4$, and of $D_n(6)\setminus A_n(6)$ for every $n\geq3$.
Thus the equality in all degrees conjectured in Satoh's 2018 report is false
at every rank in its stated range. The construction uses a classical
degree-five Lie identity, elementary free-group automorphisms, and integral
Magnus coefficients. The proof works with actual powers of $J$, rather than
only values along one-parameter representations.
\end{abstract}

\section{The conjecture and the counterexamples}
Write $F_n=\langle x_1,\ldots,x_n\rangle$ and
$\Gamma_1F_n=F_n$, $\Gamma_{r+1}F_n=[\Gamma_rF_n,F_n]$.
All lower central subgroups in this note are integral. Put
\[
 A_n(k)=\ker\bigl(\Aut F_n\longrightarrow
                  \Aut(F_n/\Gamma_{k+1}F_n)\bigr).
\]
For a representation $\rho:F_n\to\SL(2,\C)$, let
$a_{ij}(g)(\rho)=\rho(g)_{ij}$.
Let $R$ be the $\Q$-algebra of functions generated by all $a_{ij}(g)$, and set
\[
 J=(a_{ij}(g)-\delta_{ij}\mid g\in F_n,\ 1\leq i,j\leq2)\subset R,
 \qquad
 D_n(k)=\ker\bigl(\Aut F_n\longrightarrow\Aut(J/J^{k+1})\bigr).
\]
We use the substitution action $a_{ij}(g)^\phi=a_{ij}(\phi(g))$.
These are the component-function algebra and the full ideal quotient of
\cite[pp.~88--89]{SatohOWR}, not a trace algebra or a single associated
graded piece. That source records $A_n(k)\subseteq D_n(k)$ and equality for
$1\leq k\leq4$, then conjectures equality for all $k\geq1$ in the discussion
with $n\geq3$. See also \cite[\S4.2, Theorem~6.3]{SatohPreprint} for the
definitions and low-degree comparison.

Throughout, $[g,h]=ghg^{-1}h^{-1}$ for groups and $[U,V]=UV-VU$ for algebras.
Define the words
\begin{align}
 w&=[[[x_1,x_2],[x_1,x_3]],x_1],\label{eq:w}\\
 v&=[[[x_1,x_2],[x_1,[x_1,x_2]]],x_1].\label{eq:v}
\end{align}
\begin{theorem}\label{thm:main}
For every $n\geq4$, the automorphism fixing all generators except
$x_4\mapsto x_4w$ belongs to $D_n(5)\setminus A_n(5)$.
For every $n\geq3$, the automorphism fixing all generators except
$x_3\mapsto x_3v$ belongs to $D_n(6)\setminus A_n(6)$.
\end{theorem}

The degree-five identity used below is classical. With a change of sign and
variables it is precisely Burde's identity~(4); his Theorem~4.14 attributes
its occurrence in a basis of $\mathfrak{sl}_2$ identities to Razmyslov
\cite{Burde}. We report that historical attribution through the inspected
Burde source. The weighted substitution yielding the degree-six identity,
Cayley--Hamilton, leading commutator expansions and Magnus expansions are
also classical ingredients. No new polynomial identity or general PI method
is claimed here. The contribution presented is their explicit application
to the equality formulated in \cite{SatohOWR}.

\section{The matrix identity and actual ideal powers}
\begin{lemma}\label{lem:pi}
For $2\times2$ matrices over any commutative ring,
\[
 P(X,Y,Z):=[[[X,Y],[X,Z]],X]=0.
\]
Consequently $Q(X,Y):=P(X,Y,[X,Y])=0$.
\end{lemma}
\begin{proof}
Put $t=\operatorname{tr}X$, $U=[X,Y]$ and $V=[X,Z]$.
Cayley--Hamilton gives $X^2-tX+\det(X)I=0$, whence
\[
 XU+UX=[X^2,Y]=tU,\qquad XV+VX=tV.
\]
It follows that
$XUV=(tU-UX)V=tUV-U(tV-VX)=UVX$.
Likewise $XVU=VUX$, so $X$ commutes with $UV-VU$.
The second assertion is substitution. This elementary proof uses no division
or diagonalization.
\end{proof}

\begin{lemma}\label{lem:filtered}
Let $I$ be an ideal in a commutative ring $S$. If $1+u$ and $1+z$ are
invertible matrices, with $u\in\Mat_2(I^r)$ and $z\in\Mat_2(I^s)$,
$r,s\geq1$, then
\[
 [1+u,1+z]-1=(uz-zu)(1+u)^{-1}(1+z)^{-1}.
\]
Its entries lie in $I^{r+s}$, and its class modulo $I^{r+s+1}$ is the
bracket of the leading classes of $u$ and $z$.
\end{lemma}
\begin{proof}
Subtract $(1+z)(1+u)$ from $(1+u)(1+z)$ and multiply on the right by the
two displayed inverses. Both inverses are congruent to $1$ modulo $I$,
so their product does not change the class in degree $r+s$.
\end{proof}

Let $M(g)=(a_{ij}(g))\in\SL_2(R)$ be the universal matrix of a word.
Write $M(x_i)=1+B_i$, so $B_i\in\Mat_2(J)$, and let $X_i$ be its
degree-one class in $\Mat_2(\operatorname{gr}_J R)$.
Successive applications of Lemma~\ref{lem:filtered} show that the degree-five
class of $M(w)-1$ is $P(X_1,X_2,X_3)$, and the degree-six class of
$M(v)-1$ is
\[
 [[[X_1,X_2],[X_1,[X_1,X_2]]],X_1]
   =P(X_1,X_2,[X_1,X_2]).
\]
The latter substitution has weight two, giving total degree six.
The associated graded ring is commutative, so Lemma~\ref{lem:pi} applies.
The zero leading classes give the actual ideal containments
\begin{equation}\label{eq:ideal}
 M(w)-1\in\Mat_2(J^6),\qquad
 M(v)-1\in\Mat_2(J^7).
\end{equation}
No assertion that a graded ring is a polynomial ring, no smoothness or
reducedness assumption, and no passage from vanishing along arcs is needed.

\section{Integral lower central obstructions}
Let $T=\Z\langle\!\langle X_1,\ldots,X_n\rangle\!\rangle$ be the ring
of noncommutative formal power series, filtered by total degree.
The assignment $x_i\mapsto1+X_i$ defines a group homomorphism
$E:F_n\to T^\times$, since $1+X_i$ has its formal geometric inverse.
The same commutator calculation, now in the degree filtration, gives
\begin{equation}\label{eq:magnus}
 E(\Gamma_rF_n)\subseteq1+T_{\geq r}
 \qquad(r\geq1).
\end{equation}
This follows inductively because $[1+T_{\geq r},1+T_{\geq s}]
\subseteq1+T_{\geq r+s}$. Only this implication is required; neither
injectivity of $E$ nor a converse dimension-subgroup theorem is used.

The leading terms of $E(w)-1$ and $E(v)-1$ are the free associative
polynomials $P(X_1,X_2,X_3)$ and $Q(X_1,X_2)$, respectively. They do
not vanish in the free algebra. Indeed,
\begin{equation}\label{eq:coeff}
 [X_1X_1X_2X_1X_3]P=-1,\qquad
 [X_1X_1X_1X_2X_1X_2]Q=1,
\end{equation}
where $[m]$ denotes the coefficient of the monomial $m$.
For the first coefficient, in
$[[U,V],X_1]=UVX_1-VUX_1-X_1UV+X_1VU$ with
$U=[X_1,X_2]$, $V=[X_1,X_3]$, only $-X_1UV$ contributes to the
indicated monomial. For the second take
$U=X_1X_2-X_2X_1$ and
$V=X_1^2X_2-2X_1X_2X_1+X_2X_1^2$; only $X_1VU$ contributes,
using the first monomial in each of $V$ and $U$.
Thus \eqref{eq:magnus} proves
\begin{equation}\label{eq:nonmember}
 w\notin\Gamma_6F_n\quad(n\geq3),\qquad
 v\notin\Gamma_7F_n\quad(n\geq2).
\end{equation}
All coefficients are integral. Adjoining additional free generators leaves
these coefficients unchanged, so these assertions require no stability
or rational-isolator assumption.

\section{Automorphisms and the full quotient action}
\begin{proof}[Proof of Theorem~\ref{thm:main}]
If a word $u$ uses no $x_t$, the substitution $\phi_u(x_t)=x_tu$ fixing
the other generators is an automorphism. Its inverse sends $x_t$ to
$x_tu^{-1}$ and fixes all others: both substitutions fix $u$, and their
two compositions fix every free generator.

Apply this first to $(u,t)=(w,4)$, and then to $(u,t)=(v,3)$.
For the moving generator,
\[
 M(\phi_u(x_t))-M(x_t)=M(x_t)(M(u)-1).
\]
Equation~\eqref{eq:ideal} puts all its entries in $J^{k+1}$ for $k=5$
and $6$, respectively; every other generator matrix is fixed.

The entries of the finitely many matrices $M(x_i)$ generate $R$:
products give the entries of products of words, and inverse entries are
adjugate entries since all determinants are one. Therefore fixing these
algebra generators modulo $J^{k+1}$ fixes every element of $R/J^{k+1}$.
Explicitly, the difference between a product of substituted generators and
the original product is a sum of products each containing one generator
difference. Each summand lies in $J^{k+1}$. The trivial representation is
fixed by substitution, so $J$ is invariant. Hence $\phi_u$ acts identically
on the whole $J/J^{k+1}$, proving $\phi_u\in D_n(k)$.

Finally $x_t^{-1}\phi_u(x_t)=u$. The nonmembership statements
\eqref{eq:nonmember} show that $\phi_u$ fails to act identically on
$F_n/\Gamma_{k+1}F_n$. Thus $\phi_u\notin A_n(k)$, as required.
\end{proof}

For each $n\geq3$ the second construction already refutes equality in all
degrees. The first gives the earlier failure in degree five when $n\geq4$.
We make no claim about degree five at rank three, failure in every degree,
larger matrix sizes, character algebras, or the different comparison of
$A_n(k)$ with the lower central series of $\mathrm{IA}_n$.

\section{Verification, related work and status}
The accompanying archive preserves the original candidate packet and separate
independent mathematical and priority audits. Exact computations check the
matrix identities and finite ideal quotients. A distinct integral check uses
upper unitriangular matrices. In size six, assigning
\[
 x_1\mapsto I+E_{12}+E_{23}+E_{45},\quad
 x_2\mapsto I+E_{34},\quad x_3\mapsto I+E_{56}
\]
sends $w$ to $I-E_{16}$. In size seven, assigning
\[
 x_1\mapsto I+E_{12}+E_{23}+E_{34}+E_{56},\quad
 x_2\mapsto I+E_{45}+E_{67}
\]
sends $v$ to $I+E_{17}$. Here $E_{ij}$ is a matrix unit and other generators
map to $I$. Direct integral multiplication verifies both identities.
The $d$th lower central subgroup of the size-$d$ unitriangular group is
trivial, since commutators add distances above the diagonal. These give
independent finite certificates for \eqref{eq:nonmember}.
The archive's 489 colored-path evaluations recover all coefficients through
the relevant degrees. These exact checks supplement the written proof;
they are not a formal proof certificate or a literature certificate.

The bounded priority audit inspected the literal 2018 conjecture, Satoh's
public 2017 preprint version, and classical PI and related automorphism
literature. It found no verified earlier complete application to this exact
conjecture. Full journal texts of \cite{Satoh2017} and the 2021 publication
related to \cite{SatohPreprint}, as well as some related 2025--2026 journal
articles, were unavailable. Earlier versions, abstracts and related shorter
sources are identified separately in the archive. No firstness, exhaustive
novelty clearance, or continuously open status through 2026 is certified.
The theorem is a negative answer to the equality as formulated in the
inspected 2018 source.

\paragraph{AI assistance and review.}
AI tools were used extensively in developing the constructions, checking
proofs and computations, investigating sources, and writing this note.
This is an unrefereed preprint without independent human peer review.
Internal adversarial AI reviews do not constitute human peer review.
The authored note and supporting materials are released under CC BY 4.0;
third-party full-text sources are not redistributed in the archive.

\begin{thebibliography}{9}
\bibitem{SatohOWR}
T.~Satoh, On the Johnson homomorphisms of the automorphism groups of free
groups, in \emph{Topology of Arrangements and Representation Stability},
Oberwolfach Reports \textbf{15} (2018), report~2, pp.~87--90.
\href{https://doi.org/10.4171/OWR/2018/2}{doi:10.4171/OWR/2018/2}.

\bibitem{SatohPreprint}
T.~Satoh, On the graded quotients of the $\SL(m,\C)$-representation
algebras of groups, \href{https://arxiv.org/abs/1607.05411v2}{arXiv:1607.05411v2}
(8 December 2017), especially \S4.2 and Theorem~6.3.
Related journal publication: \emph{Pacific J. Math.} \textbf{312} (2021),
477--504, \href{https://doi.org/10.2140/pjm.2021.312.477}{doi:10.2140/pjm.2021.312.477}.
The cited preprint version was inspected; the full final journal text was not.

\bibitem{Satoh2017}
T.~Satoh, On the universal $\SL_2$-representation rings of free groups,
\emph{Proc. Edinburgh Math. Soc.} \textbf{60} (2017), 973--1001.
\href{https://doi.org/10.1017/S0013091516000456}{doi:10.1017/S0013091516000456}.
Bibliographic background; full text unavailable in the present audit.

\bibitem{Burde}
D.~Burde, Derivation double Lie algebras,
\href{https://arxiv.org/abs/1411.0950v2}{arXiv:1411.0950v2}
(27 February 2015), identity~(4) and Theorem~4.14.
Related journal publication:
\href{https://doi.org/10.1142/S0219498816501140}{doi:10.1142/S0219498816501140}.
\end{thebibliography}
\end{document}
